\documentclass[11pt]{article}

\usepackage[margin=1in]{geometry}
\usepackage{amsmath,amssymb,amsfonts}
\usepackage{booktabs}
\usepackage{algorithm}
\usepackage{algorithmic}
\usepackage{url}
\usepackage{hyperref}
\usepackage{setspace}
\onehalfspacing

\title{An Optimized Adaptive Heun Integrator for Differentiable\\
       Simulation of Stiff Power Electronic Systems}

\author{Hamad Alduaij\\
School of Electrical, Computer and Energy Engineering\\
Arizona State University, Tempe, AZ\\
\texttt{halduaij@asu.edu}}

\date{}

\begin{document}

\maketitle

\begin{abstract}
Differentiable ODE solvers are central to gradient-based training of
neural controllers for power electronic systems.  We identify a
first-same-as-last (FSAL) implementation error in the widely used
\texttt{torchdiffeq} adaptive Heun solver that evaluates the first
stage at an incorrect point, inflating the error constant by
2--3$\times$.  We provide a corrected and optimized implementation
with JIT-compiled step logic, branchless adaptive step-size control,
gradient separation, and Horner-form dense output.  On an 18-state
system with stiff phase-locked-loop dynamics, the optimized solver
achieves a 2.1$\times$ wall-clock speedup over the original
\texttt{torchdiffeq} implementation while producing 2--3$\times$
smaller errors at the same step size.  The code is open-source at
\url{https://github.com/halduaij/inverter-control}.
\end{abstract}

% ============================================================================
\section{Introduction}
\label{sec:intro}
% ============================================================================

Gradient-based optimization of neural controllers for power systems
requires differentiable time-domain simulation.  The controller
parameters appear inside the ODE right-hand side $f(t, y; \theta)$, and
training proceeds by backpropagating a trajectory-level loss through the
ODE integrator to obtain $\partial \mathcal{L}/\partial\theta$.
Libraries such as \texttt{torchdiffeq}~\cite{Chen18NeuralODE} provide
adaptive-step ODE solvers within the PyTorch autograd framework and have
become standard infrastructure for neural ODE training.

Power electronic systems, however, present a specific numerical
challenge: the coexistence of fast inner control loops (phase-locked
loops with bandwidth $\sim$900~rad/s, current controllers at
$\sim$1~kHz) and slow electromechanical dynamics (governor response at
$\sim$0.1--1~Hz).  This moderate stiffness makes explicit methods
viable---unlike the deeply stiff chemical kinetics problems that mandate
implicit solvers---but demands that the explicit method be both correct
and efficient.  The generic Runge--Kutta infrastructure in
\texttt{torchdiffeq} introduces per-step overhead (Butcher tableau
iteration, Python-level branching for step control, redundant tensor
allocation) that compounds over the thousands of steps needed for
moderately stiff simulations.

In this note we report two findings regarding \texttt{torchdiffeq}'s
adaptive Heun solver.  First, we identify a FSAL (first-same-as-last)
implementation error that evaluates $k_1$ at the wrong point, inflating
the per-step error constant by a factor of 2--3$\times$ while preserving
second-order convergence.  Second, and more significant in practice, we
provide an optimized standalone implementation that eliminates the
generic RK overhead through JIT-compiled step arithmetic, branchless
adaptive control, and gradient separation.

The contributions are:
\begin{enumerate}
\item Identification and correction of the FSAL evaluation-point error
      in the \texttt{torchdiffeq} Heun solver (Section~\ref{sec:bug}).
\item A JIT-compiled adaptive Heun solver with branchless step-size
      control, gradient separation, and Horner-form dense output that
      achieves 2.1$\times$ wall-clock speedup over \texttt{torchdiffeq}
      (Section~\ref{sec:impl}).
\item Validation on an 18-state system with stiff PLL dynamics
      (Section~\ref{sec:results}).
\end{enumerate}

% ============================================================================
\section{Background: Adaptive Heun Method}
\label{sec:background}
% ============================================================================

Heun's method is the simplest second-order explicit Runge--Kutta scheme.
Given an ODE $\dot{y} = f(t, y)$, one step from $(t_n, y_n)$ with
step size $h$ proceeds as:
\begin{align}
  k_1 &= f(t_n,\; y_n), \label{eq:k1} \\
  k_2 &= f(t_n + h,\; y_n + h\,k_1), \label{eq:k2} \\
  y_{n+1} &= y_n + \tfrac{h}{2}(k_1 + k_2). \label{eq:heun}
\end{align}
The Butcher tableau is
\[
\begin{array}{c|cc}
0 & 0 & 0 \\
1 & 1 & 0 \\
\hline
  & \tfrac12 & \tfrac12
\end{array}
\]
with the embedded first-order (Euler) method $\hat{y}_{n+1} = y_n + h\,k_1$
providing a local error estimate
\begin{equation}
  e_{n+1} = y_{n+1} - \hat{y}_{n+1} = \tfrac{h}{2}(k_2 - k_1).
  \label{eq:err}
\end{equation}

\subsection{FSAL Property}

For some Runge--Kutta methods (e.g., Dormand--Prince), the last stage
evaluation $k_s$ is computed at the accepted solution $y_{n+1}$, so it
can be reused as $k_1$ for the next step.  This ``first same as last''
(FSAL) property saves one function evaluation per accepted step.

For Heun's method, $k_2$ is evaluated at the Euler predictor
$\tilde{y} = y_n + h\,k_1$, which differs from the accepted solution
$y_{n+1}$ by $O(h^2)$.  Reusing $k_2$ as $k_1$ for the next step
therefore evaluates the first stage at the wrong point, introducing an
$O(h^2)$ per-step error in $k_1$.  Since $k_1$ is multiplied by $h$ in
the update formula, this contributes an $O(h^3)$ local truncation error
per step---the same order as Heun's inherent truncation error---so the
global convergence order remains 2.  However, the additional error
increases the leading error constant by a factor of 2--3$\times$,
producing less accurate solutions at the same step size.

% ============================================================================
\section{The FSAL Bug in torchdiffeq}
\label{sec:bug}
% ============================================================================

The \texttt{torchdiffeq} library implements adaptive Heun as a subclass
of its generic adaptive-step Runge--Kutta solver.  The generic solver
assumes the FSAL property holds and caches the last stage evaluation for
reuse.  For methods where FSAL is valid (Dormand--Prince 4/5), this is
correct.  For Heun's method, however, the cached $k_2$ was evaluated
at $\tilde{y} = y_n + h\,k_1$, not at $y_{n+1}$.  Reusing it as
$k_1$ for the next step means
\begin{equation}
  k_1^{(n+1)} = f(t_{n+1},\; \tilde{y}_n) \neq f(t_{n+1},\; y_{n+1}),
  \label{eq:fsal_bug}
\end{equation}
where $\tilde{y}_n - y_{n+1} = \tfrac{h}{2}(k_1 - k_2) = O(h^2)$.
This shifts the evaluation point of $k_1$ by a second-order amount
at each step.  While this does not degrade the convergence order
(both versions remain second-order), it inflates the error constant,
producing 2--3$\times$ larger errors at any given step size
(Table~\ref{tab:order}).  The corrupted $k_1$ also affects the error
estimator~\eqref{eq:err}, which uses $k_1$ in the difference
$k_2 - k_1$, potentially reducing the reliability of adaptive step-size
selection.

The fix is straightforward: evaluate $k_1 = f(t_{n+1}, y_{n+1})$ fresh
at each step, at the cost of one additional function evaluation per
step (2 evaluations total, matching the theoretical cost of a
two-stage method).

% ============================================================================
\section{Implementation}
\label{sec:impl}
% ============================================================================

Beyond the FSAL correction, we provide a standalone adaptive Heun solver
that replaces the generic Runge--Kutta infrastructure with specialized
routines.  The \texttt{torchdiffeq} generic solver iterates over the
Butcher tableau, allocates intermediate $k$-tensors, and uses
Python-level branching for step acceptance---all unnecessary overhead for
a two-stage method.  The following optimizations eliminate this overhead.

\subsection{JIT-Compiled Step Logic}

All arithmetic in the step function---Euler prediction, Heun correction,
error estimation, and step-size update---is decorated with
\texttt{@torch.jit.script}, eliminating Python interpreter overhead for
the inner loop.  The Heun step is computed directly as
$y_1 = y_0 + \tfrac{h}{2}(k_1 + k_2)$ without iterating over a
Butcher tableau.  Optionally, the entire step can be wrapped with
\texttt{torch.compile} (PyTorch 2.0+) for further kernel fusion.

\subsection{Branchless Step-Size Control}

The step-size update uses \texttt{torch.where} throughout to avoid
Python-level branching that would cause CPU--GPU synchronization:
\begin{equation}
  h_{n+1} = h_n \cdot \text{clamp}\!\bigl(
    \sigma \cdot r_n^{-1/p},\; d_{\min},\; d_{\max}
  \bigr),
  \label{eq:step_ctrl}
\end{equation}
where $r_n = \|e_{n+1}\|_{\text{RMS}} / \text{tol}$ is the normalized
error ratio, $p = 2$ is the method order, $\sigma = 0.9$ is the safety
factor, $d_{\min} = 0.2$, and $d_{\max} = 10$.  The tolerance is
mixed absolute--relative:
$\text{tol}_i = \text{atol} + \text{rtol} \cdot
\max(|y_n^{(i)}|, |y_{n+1}^{(i)}|)$.

A step is accepted if $r_n \le 1$ or $h_n$ is already at the minimum
allowed step size.

\subsection{Dense Output}

For saving solutions at prescribed output times without restricting the
adaptive step size, we fit a quartic Hermite interpolant through
$y_n$, $y_{n+1}$, $y_{\text{mid}} = y_n + (h/2)\,k_1$, $k_1$, and
$k_2$, evaluated via Horner's rule with \texttt{torch.addcmul} for
fused multiply-add:
\begin{equation}
  y(\tau) = c_0 + \tau\bigl(c_1 + \tau(c_2 + \tau(c_3 + \tau\,c_4))\bigr),
  \quad \tau = \frac{t - t_n}{h} \in [0, 1].
  \label{eq:horner}
\end{equation}

\subsection{Gradient Separation}

The accept/reject decision~\eqref{eq:step_ctrl} is non-differentiable.
We run the control logic inside \texttt{torch.no\_grad()} for
efficiency, then recompute the accepted solution $y_{n+1} = y_n +
(h/2)(k_1 + k_2)$ outside the no-grad block so that gradients flow
through the Heun formula to $f$ and ultimately to the controller
parameters $\theta$.  Algorithm~\ref{alg:step} summarizes one adaptive
step.

\begin{algorithm}[t]
\caption{Corrected Adaptive Heun Step}
\label{alg:step}
\begin{algorithmic}[1]
\REQUIRE State $y_n$, time $t_n$, step size $h$, tolerances atol, rtol
\STATE $k_1 \gets f(t_n, y_n)$ \COMMENT{Fresh evaluation (no FSAL reuse)}
\STATE $\tilde{y} \gets y_n + h\,k_1$ \COMMENT{Euler predictor}
\STATE $k_2 \gets f(t_n + h,\; \tilde{y})$
\STATE \textbf{with} \texttt{no\_grad()}:
\STATE \quad $y_1 \gets y_n + \tfrac{h}{2}(k_1 + k_2)$
\STATE \quad $e \gets \tfrac{h}{2}(k_2 - k_1)$
\STATE \quad $r \gets \text{RMS}(e\,/\,\text{tol})$
\STATE \quad $h_{\text{next}} \gets h \cdot
              \text{clamp}(0.9\,r^{-0.5},\; 0.2,\; 10)$
\STATE \quad accept $\gets (r \le 1)$
\STATE \textbf{end}
\STATE $y_1^{\text{grad}} \gets y_n + \tfrac{h}{2}(k_1 + k_2)$
       \COMMENT{With gradients}
\RETURN $y_1^{\text{grad}}$, accept, $h_{\text{next}}$
\end{algorithmic}
\end{algorithm}

% ============================================================================
\section{Application: Multi-Converter Power System}
\label{sec:application}
% ============================================================================

The target application is a multi-converter power system with
distributed adaptive voltage-oriented control (dVOC).  Each converter
has voltage and current PI control loops, phase estimation, and filter
dynamics.  We benchmark on an 18-state system that captures the
essential stiffness structure: six fast PLL states with eigenvalues
near $-30 \pm 780j$~rad/s, six medium-speed filter states, and six
slow power-sharing states at 1--10~rad/s, producing a stiffness ratio
of $\sim$100--1000.

This stiffness ratio is mild enough that explicit methods remain
competitive with implicit solvers (which would require Jacobian
factorization at each step), but severe enough that solver overhead
matters---thousands of steps are needed for a 100~ms simulation at
tight tolerances.  The per-step overhead of generic Runge--Kutta
infrastructure (tableau iteration, tensor allocation, Python branching)
accumulates to a measurable fraction of total wall-clock time.

The solver is embedded in a differentiable training loop where neural
network gain schedulers are optimized via backpropagation through the
ODE trajectory.  Both solution accuracy and computational efficiency
directly affect training: larger error constants produce noisier
gradients, while per-step overhead limits the number of training
iterations per hour.

% ============================================================================
\section{Numerical Results}
\label{sec:results}
% ============================================================================

\subsection{Convergence Order and Error Constants}

We verify the convergence order and error constants on three test
problems with known solutions, integrated with fixed step sizes.
Table~\ref{tab:order} reports results for $\dot{y} = -0.5\,y$,
$y(0) = 1$, integrated to $t = 4$ (exact solution $y(4) = e^{-2}$).

\begin{table}[t]
\centering
\caption{Global error vs.\ step size: original (FSAL reuse) vs.\
corrected.  Both methods converge at order 2 ($\sim$4$\times$ error
reduction per halving of $h$); the FSAL bug inflates the error constant
by 2.2--2.5$\times$.}
\label{tab:order}
\small
\begin{tabular}{rccc}
\toprule
$h$ & Original (FSAL) & Corrected & Error ratio \\
\midrule
0.400 & $4.7 \times 10^{-3}$ & $2.1 \times 10^{-3}$ & 2.2$\times$ \\
0.200 & $1.2 \times 10^{-3}$ & $4.9 \times 10^{-4}$ & 2.4$\times$ \\
0.100 & $2.9 \times 10^{-4}$ & $1.2 \times 10^{-4}$ & 2.4$\times$ \\
0.050 & $7.1 \times 10^{-5}$ & $2.9 \times 10^{-5}$ & 2.5$\times$ \\
0.025 & $1.8 \times 10^{-5}$ & $7.1 \times 10^{-6}$ & 2.5$\times$ \\
\bottomrule
\end{tabular}
\end{table}

Both versions exhibit the expected $\sim$4$\times$ error reduction when
$h$ is halved, confirming second-order convergence.  The error ratio
(FSAL to corrected) is consistently 2.2--2.5$\times$ across all step
sizes, indicating that the FSAL bug inflates the leading error constant
without changing the convergence order.  We observe similar ratios on
nonlinear ($\dot{y} = -y^2$, ratio 2.9--3.0$\times$) and oscillatory
($\ddot{y} + y = 0$, ratio 2.5$\times$) test problems.

\subsection{Performance vs.\ torchdiffeq}

Table~\ref{tab:perf} compares wall-clock time for the optimized solver
against \texttt{torchdiffeq}'s adaptive Heun on the 18-state stiff
system (100~ms integration, 50 output points, CPU).  The optimized
solver achieves a consistent 2.1$\times$ speedup despite performing
2 function evaluations per step (vs.\ 1 for the FSAL-reusing original).
The speedup comes from eliminating generic RK overhead: specialized
two-stage step logic replaces Butcher tableau iteration, JIT-compiled
arithmetic removes Python interpreter overhead, and branchless step-size
control avoids CPU--GPU synchronization points.

\begin{table}[t]
\centering
\caption{Wall-clock time: \texttt{torchdiffeq} vs.\ optimized solver
(18-state system, 100~ms integration, CPU).}
\label{tab:perf}
\small
\begin{tabular}{lrrr}
\toprule
\textbf{Solver} & \textbf{Time (ms)} & \textbf{Speedup} & \textbf{Steps} \\
\midrule
\texttt{torchdiffeq} (rtol=$10^{-4}$) & 755 & 1.0$\times$ & $\sim$4300 \\
Optimized (rtol=$10^{-4}$)             & 350 & 2.2$\times$ & 4290 \\
\midrule
\texttt{torchdiffeq} (rtol=$10^{-6}$) & 7328 & 1.0$\times$ & $\sim$43{,}600 \\
Optimized (rtol=$10^{-6}$)             & 3441 & 2.1$\times$ & 43{,}639 \\
\bottomrule
\end{tabular}
\end{table}

Both solvers take approximately the same number of steps for a given
tolerance, confirming that the speedup is due to reduced per-step
overhead rather than algorithmic step count differences.

% ============================================================================
\section{Conclusion}
\label{sec:conclusion}
% ============================================================================

We identified a FSAL implementation error in the \texttt{torchdiffeq}
adaptive Heun solver that evaluates the first stage at the Euler
predictor rather than the accepted solution, inflating the error
constant by 2--3$\times$.  The corrected and optimized implementation,
available at \url{https://github.com/halduaij/inverter-control},
combines the FSAL fix with JIT-compiled step logic, branchless
adaptive control, and gradient separation.  On an 18-state system
with stiff PLL dynamics, the optimized solver achieves 2.1$\times$
wall-clock speedup over \texttt{torchdiffeq} while producing more
accurate solutions, directly improving both simulation throughput and
gradient quality for neural controller training.

% ============================================================================
\bibliographystyle{IEEEtran}

\begin{thebibliography}{1}

\bibitem{Chen18NeuralODE}
R.~T.~Q. Chen, Y.~Rubanova, J.~Bettencourt, and D.~Duvenaud,
``Neural ordinary differential equations,'' in
\textit{Advances in Neural Information Processing Systems (NeurIPS)},
vol.~31, 2018.

\bibitem{Paszke19PyTorch}
A.~Paszke \textit{et al.},
``{PyTorch}: An imperative style, high-performance deep learning library,'' in
\textit{Advances in Neural Information Processing Systems (NeurIPS)},
vol.~32, 2019, pp.~8024--8035.

\bibitem{DormandPrince80}
J.~R. Dormand and P.~J. Prince,
``A family of embedded {Runge--Kutta} formulae,''
\textit{J. Comput. Appl. Math.}, vol.~6, no.~1, pp.~19--26, 1980.

\bibitem{HairerNorsettWanner93}
E.~Hairer, S.~P. N{\o}rsett, and G.~Wanner,
\textit{Solving Ordinary Differential Equations I: Nonstiff Problems},
2nd~ed.\hspace{0.5em} Berlin: Springer, 1993.

\end{thebibliography}

\end{document}
